\documentclass[10pt]{article}
\usepackage[T1]{fontenc}
\usepackage[utf8]{inputenc}
\usepackage[french]{babel}
\usepackage{amsmath,amssymb,mathrsfs,mathpazo}
\usepackage{geometry}
\geometry{a4paper,margin=2cm}
\usepackage{graphicx,array,tabularx,tabularray,colortbl,multirow,multicol,booktabs,enumitem,pifont,ulem,fancybox,diagbox,cancel}
\usepackage{pstricks,pstricks-add,pst-plot,pst-tree,pst-node,pst-text,pst-func,pst-math,pst-bspline,pst-3dplot,pst-eucl,pst-solides3d}
\let\pscalculate\undefined
\usepackage{pst-bezier}
\usepackage{tikz,pgfplots,tkz-tab}
\usetikzlibrary{arrows,arrows.meta,patterns,calc,positioning,trees,decorations.markings,shapes.geometric,shapes.misc,shapes.arrows,intersections,decorations.pathmorphing,decorations.pathreplacing,backgrounds,shadows,babel}
\pgfplotsset{compat=1.18}
\usepgfplotslibrary{fillbetween}
\usepackage[np]{numprint}
\usepackage{listings,makecell,pas-tableur,twoopt}
\let\barre\undefined
\usepackage[pstricks]{tablvar}
\definecolor{darkgreen}{RGB}{0,100,0}
\providecommand{\red}{\color{red}}
\providecommand{\blue}{\color{blue}}
\providecommand{\green}{\color{green}}
\providecommand{\black}{\color{black}}
\providecommand{\white}{\color{white}}
\providecommand{\gray}{\color{gray}}
\providecommand{\cyan}{\color{cyan}}
\providecommand{\magenta}{\color{magenta}}
\providecommand{\yellow}{\color{yellow}}
\providecommand{\violet}{\color{violet}}
\providecommand{\orange}{\color{orange}}
\providecommand{\euro}{\texteuro}
\providecommand{\EUR}{\texteuro}
\providecommand{\ieme}{\textsuperscript{e}}
\providecommand{\degres}{\ensuremath{^\circ}}
\providecommand{\degree}{\ensuremath{^\circ}}
\providecommand{\vv}[1]{\overrightarrow{#1}}
\providecommand{\Coord}[1]{#1}
\providecommand{\decosix}{}
\providecommand{\starredbullet}{\ensuremath{\bullet}}
\pagestyle{empty}
\setlength{\parindent}{0pt}
\usepackage{adjustbox}\DeclareUnicodeCharacter{2236}{:}\DeclareUnicodeCharacter{2190}{\ensuremath{\leftarrow}}
\providecommand{\R}{\mathbb{R}}
\providecommand{\N}{\mathbb{N}}
\providecommand{\Z}{\mathbb{Z}}
\providecommand{\Q}{\mathbb{Q}}
\providecommand{\D}{\mathbb{D}}
\providecommand{\C}{\mathbb{C}}
\providecommand{\e}{{\mathrm{e}}}
\providecommand{\vect}[1]{\overrightarrow{#1}}
\providecommand{\vectt}[1]{\overrightarrow{\mathrm{#1}}}
\providecommand{\vvt}[1]{\overrightarrow{\mathrm{#1}}}
\providecommand{\np}[2][]{{#2~#1}}
\providecommand{\vv}[1]{{\overrightarrow{#1}}}
\renewcommand{\d}{{\mathrm{d}}}
\providecommand{\degres}{{^\circ}}
\providecommand{\degre}{{^\circ}}
\providecommand{\blue}{\color{blue}}
\providecommand{\red}{\color{red}}
\providecommand{\enskip}{\quad}
\providecommand{\llbracket}{\lbrack\!\lbrack}
\providecommand{\rrbracket}{\rbrack\!\rbrack}
\providecommand{\euro}{{\text{€}}}
\providecommand{\no}{{\text{n°}}}
\providecommand{\interval}[2]{{[#1;#2]}}
\providecommand{\intervFO}[2]{{[#1;#2[}}
\providecommand{\dotfill}{\ldots}

\providecommand{\ds}{\displaystyle}
\providecommand{\textding}[1]{#1}
\providecommand{\Oij}{\(\mathrm{O};\vec{\imath},\vec{\jmath}\)}
\providecommand{\Oijk}{\(\mathrm{O};\vec{\imath},\vec{\jmath},\vec{k}\)}
\providecommand{\Ouv}{\(\mathrm{O};\vec{u},\vec{v}\)}
\providecommand{\intervFF}[2]{\ensuremath{[#1;#2]}}
\renewcommand{\euro}{\ifmmode\text{\texteuro}\else\texteuro\fi}
\DeclareUnicodeCharacter{00B0}{\ifmmode^{\circ}\else\textdegree\fi}
\DeclareUnicodeCharacter{20AC}{\ifmmode\text{\texteuro}\else\texteuro\fi}
% Énoncés : Éducation nationale. Transcription et corrigés : APMEP, Denis Vergès et contributeurs. Découpage et classement : CoopMaths, Sébastien Lozano, Éric Elter et contributeurs. Adaptation de la présentation : Maths Explorer.
% CC BY-SA : https://coopmaths.fr/www/about/
% https://coopmaths.fr/alea/static/bac/2025/tex/bac_2025_05_sujet2_ameriquenord_2_cor.tex
\begin{document}
\noindent\textit{APMEP / CoopMaths. CC BY-SA. Présentation adaptée par Maths Explorer.}\par\medskip
On désigne par $\left(u_{n}\right)$ la suite définie par :

\[u_{0}=2 \quad \text { et, pour tout entier naturel } n,\quad  u_{n+1}=\sqrt{u_{n}}\]

\medskip

\textbf{Partie A}

\medskip

\begin{enumerate}
\item
	\begin{enumerate}
		\item 

$u_1=\sqrt{u_0}=\sqrt{2}$

$u_2=\sqrt{u_1} =\sqrt{\sqrt{2}}$.
		\item À l'aide de la calculatrice, on peut conjecturer que la suite est décroissante et converge vers 1.
	\end{enumerate}
\item
	\begin{enumerate}
		\item Posons, pour tout entier naturel $n$, l'affirmation $P_n$ :~\og $1 \leqslant u_{n+1} \leqslant u_{n}$\fg{}.\hfill

\emph{Initialisation :} On a calculé $u_1 = \sqrt{2} \approx 1,414$ et $u_0=2$

On a donc bien $1 \leqslant u_1 \leqslant u_0$.

L'affirmation $P_0$ est donc vraie.

\emph{Hérédité :} Soit $n$ un entier naturel non nul. On suppose que l'affirmation $P_n$ est vraie, c'est-à-dire que $1 \leqslant u_{n+1} \leqslant u_{n}$.

Montrons que l'affirmation $P_{n+1}$ est vraie, c'est à dire que $1 \leqslant u_{n+2} \leqslant u_{n+1}$

Par hypothèse de récurrence, on a :

$ \begin{aligned}[t]
1 \leqslant u_{n+1} \leqslant u_n
	&\implies \sqrt{1} \leqslant \sqrt{u_{n+1}} \leqslant \sqrt{u_n}\\
	&\qquad \text{car  la fonction racine carrée est croissante sur } \R^+ \\
	&\implies 1 \leqslant u_{n+2} \leqslant u_{n+1}
\end{aligned}$

Si, pour $n$ naturel, $P_n$ est vraie, alors $P_{n+1}$ est vraie également.

\emph{Conclusion :} L'affirmation $P_0$ est vraie, et, pour tout entier naturel non nul $n$, la véracité de l'affirmation est héréditaire : d'après le principe de récurrence, on en déduit que pour tout entier naturel non nul $n$, on a : $1 \leqslant u_{n+1} \leqslant u_{n}$.
	\item D'après la question précédente, pour tout entier $n$, $1\leqslant u_n$ donc la suite $\left(u_n\right)$ est minorée par 1.

Et pour tout entier $n$, $u_{n+1}\leqslant u_n$ donc la suite $\left( u_n\right)$ est décroissante.

Or toute suite décroissante et minorée est convergente donc, la suite $\left( u_n\right)$ est convergente vers une limite $\ell \geqslant 1$.
	\item Soit $x \in [0~;~+\infty[$ :

$\begin{aligned}\sqrt{x} =x &\iff x =x^2 \text { car }x\geqslant 0 \\
	&\iff x^2-x=0\\
	& \iff x(x-1)=0\\
	&\iff x=0 \text{ ou } x=1
\end{aligned}$.

L'équation $\sqrt{x}=x$ admet deux solutions sur $[0~;~+\infty[$ : 0 et 1.
	\item La suite $\left(u_{n}\right)$ est convergente vers $\ell$ et est définie par récurrence par :

$\forall n \in \N, \quad u_{n+1}=f(u_n)$ \quad où  $f$ est la fonction racine carrée qui est continue sur $[0~;~+\infty[$.

D'après le théorème du point fixe, $\ell$ est solution de l'équation $f(x) = x$.

Nous avons résolu cette équation à la question précédente, les solutions sont 0 et 1, or d'après la question \textbf{A. 1. b.}, $\ell \geqslant 1$ donc $\ell =1$.

donc : $\lim\limits_{n\to +\infty} u_n = 1$
	\end{enumerate}
\end{enumerate}

\medskip

\textbf{Partie B}

\medskip

\begin{enumerate}
\item
	\begin{enumerate}
		\item Soit $n \in \N$ : $\begin{aligned}[t] v_{n+1}& =\ln \left(u_{n+1}\right)\\
	&= \ln \left( \sqrt{u_{n}}\right)\\
	& = \dfrac{1}{2}\ln\left(u_n\right)\\
	&= \dfrac{1}{2}v_n
\end{aligned}$

La suite $\left(v_{n}\right)$ est donc géométrique de raison $q=\dfrac{1}{2}$ et de premier terme \linebreak $v_0=\ln\left(u_0\right)=\ln(2)$.
		\item On a donc, pour tout entier naturel $n$, $v_n=v_0\times q^n=\ln(2)\left(\dfrac{1}{2}\right)^n=\dfrac{\ln(2)}{2^n}$.
		\item  Or, pour tout entier naturel $n$,  $v_n=\ln \left(u_{n}\right)$ donc $\dfrac{\ln\left( 2 \right)}{2^n}= \ln \left(u_{n}\right)$ c'est à dire : $\ln(2)=2^n\ln\left(u_n\right)$.
	\end{enumerate}
\item
	\begin{enumerate}
		\item À l'aide de la calculatrice on trouve que pour tout $k<7$, $u_k > 1,01$ et \linebreak $u_{7}\approx \np{1,00543}$ donc $k=7$.
		\item On décide de prendre $x - 1$ comme approximation de $\ln (x)$ lorsque $x$ appartient à l'intervalle $]0,99~;~1,01[$ donc $\ln\left(u_{7}\right)\approx u_{7} - 1 \approx \np{0,00543}$.
\item D'après les questions \textbf{1. c.} et \textbf{2. b.} de la \textbf{partie B} une approximation de $\ln (2)$ est $2^7\ln\left(u_7\right)$ soit $2^7 \times \np{0,00543}$ c'est à dire $\np{0,695}$.
	\end{enumerate}
\item On généralise la méthode précédente à tout réel $a$ strictement supérieur à 1.

\begin{center}
\begin{ttfamily}
\begin{tabularx}{8cm}{|X|}\hline
\textbf{from} math \textbf{import}*\\
\textbf{def} Briggs(a):\\
\quad n = 0\\
\quad \textbf{while} a >= 1.01:\\
\quad \quad a = \textbf{sqrt}(a)\\
\quad \quad n = n+1\\
\quad L =2**n *(a-1) \\
\textbf{return} L\\ \hline
\end{tabularx}
\end{ttfamily}
\end{center}
\end{enumerate}

\bigskip

\end{document}
